\section{Discussion}
\label{sec:discussion}

\subsection{The Death of Glitch Tokens}

Our results show that the \glitchtoken phenomenon documented by \citet{rumbelow2023solidgoldmagikarp} and subsequently studied in depth \citep{land2024fishing, li2024glitch, zhang2024glitchprober, wu2024glitchminer} is entirely resolved in the GPT-4.1 model family. Every one of the 50 tested \glitchtokens---including \magikarp, the token that launched the subfield---is reproduced perfectly by all three models. This likely reflects improvements in tokenizer training, vocabulary curation, and post-training alignment over the past two years. While this is a positive development for model reliability, it also means that \glitchtoken benchmarks are no longer useful for evaluating current-generation models and that the research community should shift attention to the new failure modes we identify.

\subsection{Why Does Clean Text Fail Where Misspelled Text Succeeds?}

The most surprising finding is \gptfull's paradoxical behavior: it truncates clean varied text to $\sim$50 characters at lengths $\geq$8k while perfectly reproducing 25,000 characters of misspelled text. We consider three possible explanations:

\para{Anti-copying filters.}
The model or its serving infrastructure may apply output filters designed to prevent verbatim reproduction of copyrighted or memorized text. Clean, well-formed text is more likely to match training data, triggering these filters. Misspelled text, being clearly ``original,'' may bypass them. This explanation is consistent with the growing emphasis on copyright safety in commercial LLMs.

\para{Deduplication mechanisms.}
A related possibility is that the model detects high overlap between the input and its training data and applies a truncation policy. Misspelled text, being out-of-distribution \wrt training data, would not trigger this mechanism.

\para{Output length policies.}
The truncation may reflect API-level output length restrictions that are content-dependent. The consistent length ratio of $\approx 0.005$ across different clean text lengths suggests a fixed truncation threshold rather than a gradual degradation.

We cannot definitively distinguish these explanations without access to the model's internal architecture and serving infrastructure. However, the consistency and specificity of the behavior---perfectly reproducing misspelled text of equal or greater length---argues against a simple capability limitation.

\subsection{Repetition as a Failure Trigger}

The failure of all models on long repetitive text reveals a separate vulnerability. When input text contains many repetitions of the same passage, models exhibit two failure modes: progressive truncation (\gptomini) and hallucinated extension (\gptmini, producing 5.15$\times$ the input). These failures likely arise from the interaction between the model's attention mechanism and repeated content. In a Transformer architecture, attending over many identical segments may cause the model to lose positional grounding, leading it to either stop prematurely or fail to recognize when to stop.

This finding has practical implications for applications involving templated or repetitive content, such as form generation, data entry, and code with repeated patterns.

\subsection{Limitations}

\para{Model coverage.}
We tested only OpenAI models from the GPT-4.1 family. The clean-versus-misspelled asymmetry may be specific to OpenAI's serving infrastructure rather than a general property of large language models. Testing Claude, Gemini, and open-weight models such as Llama 3 would clarify the generality of our findings.

\para{API versus model behavior.}
The extreme-length truncation in \gptfull may reflect API-level output limits or safety filters rather than model-level capability limits. We set max\_tokens $= 16{,}384$, which should be sufficient for all tested lengths, but internal mechanisms may impose additional constraints that are not documented.

\para{Prompt sensitivity.}
We tested three prompt templates but the space of possible instructions is vast. Different phrasing---such as framing the task as code execution or using chain-of-thought---might elicit different behavior, particularly at extreme lengths.

\para{Reproducibility.}
All experiments used temperature $= 0$ for determinism. However, OpenAI's API does not guarantee bit-exact reproducibility across time even at temperature $= 0$, so our results represent a snapshot of model behavior at a specific point in time.

\para{No open-weight baselines.}
We could not test the original models (GPT-2, Llama-2-7B) where \glitchtokens were first discovered, limiting our ability to provide a direct before-and-after comparison within the same model family.

\subsection{Broader Implications}

\para{For practitioners.}
Users relying on LLMs for faithful text copying should be aware of length-dependent failure modes. Our results suggest that breaking long texts into chunks of $<$8k characters can mitigate truncation risks. The counterintuitive finding that misspelled text is reproduced more reliably than clean text is relevant for applications involving noisy or user-generated content.

\para{For model developers.}
The clean-versus-misspelled asymmetry suggests that content-dependent output policies may have unintended side effects. A policy designed to prevent copyright-infringing reproduction inadvertently makes the model less reliable for legitimate text-copying tasks.

\para{For security researchers.}
The \glitchtoken attack surface has been substantially reduced in modern models. However, the length-based and repetition-based failure modes we identify could potentially be exploited in adversarial settings, for example by constructing inputs that trigger hallucinated extension to exhaust output token budgets.
